% =====================================================================
%  Chapter 11 — Capacity: Why One Investor Is an Advantage
% =====================================================================

\chapter{Capacity: Why One Investor Is an Advantage}
\label{ch:capacity}

Chapter 10 ended with the costs that apply to this reader. This chapter
states the result that follows from their absence: a single investor
occupies a position in the cost structure that no institution\idx{Institution} can occupy,
and the advantage is large enough to change which strategies are worth
running at all.

\section{The institutional cost structure}

Frazzini, Israel and Moskowitz's study of trading costs in factor
strategies assembled roughly a trillion dollars of live data across
nineteen markets from 1998 to 2011 and compared the realised trading costs
of factor portfolios against the prior literature's estimates.

Their headline result is that actual trading costs are less than a tenth
of prior estimates. The prior estimates had been built from the visible
market impact\idx{Market impact} of large trades; the live data from institutional portfolios
showed much smaller realised costs, because the portfolios traded far less
than the estimates assumed and because much of the rebalancing could be
done internally.

Their second result is a ranking that matters enormously for what a single
investor should run. The long-horizon value and momentum premia were
highly scalable --- they can be traded in size\idx{Size} without materially
degrading the return. The short-horizon reversal premia were the most
capacity\idx{Capacity}-constrained: their returns fall fastest as capital is added,
because the trade itself is the mechanism of the effect and it must be
completed before the prices revert. The size factor sat in between, and
the bet against beta was the most scalable of all.

\section{Why the ranking reverses for one investor}

For a single investor, none of this matters, and the reason is simply
scale. A portfolio of fifty thousand dollars in a fund with a hundred
million dollars of daily volume is 0.05 per cent of one day's trading.
There is no market impact, there is no crowding, and there is no partial
execution.

But the interesting consequence is not that the single investor escapes
the costs. It is that the single investor can run strategies that the
institutional cost structure has made unviable, and can run them on a
scale where the friction is simply the spread.

Specifically:

\begin{itemize}
\item \textbf{Short-horizon signals become available.} The reversal
effect that institutions cannot scale is available at a size where the
spread cost is two basis points and the signal horizon is a week. This is
not an opportunity the book takes up --- the book's horizon is monthly ---
but it is the reason the single investor's constraint set is genuinely
different from the institution's.
\item \textbf{Implementation is a rounding error.} For an institution, the
difference between a good and a bad implementation is worth more than the
difference between a good and a bad model. For a single investor, the
implementation cost is a fraction of a basis point a year and the model is
everything. That is why this book spends five parts on the model.
\item \textbf{The decay still applies.} McLean and Pontiff's decay finding
splits into a cost component and a non-cost component. The single investor
escapes the cost component and is fully exposed to the rest. A signal that
has decayed because its participants crowded it is not re-energised by
being small.
\end{itemize}

\section{The arithmetic of the position}

It is worth being concrete about the size of the advantage, because it is
the single most under-appreciated fact in the whole subject.

Suppose a strategy has a gross expected excess return of \(\delta\) a year
and a cost of \(c\) a year, where \(c\) is proportional to the square root
of the capital deployed divided by the fund's average daily volume. For a
fund with average daily volume \(V\) and a position of size \(A\), the
square-root law gives an impact cost of order
\begin{equation}
c \approx \frac{\sigma_{\mathrm{day}}}{\sqrt{V}}\,
              \frac{A}{V}\,\eta,
\end{equation}
with \(\eta\) a constant of order one and \(\sigma_{\mathrm{day}}\) the
daily volatility.

For a broad-market ETF with \(V\) of about five billion dollars and
\(\sigma_{\mathrm{day}}\) of about 0.8 per cent, and a rebalancing
requirement of turning over the whole portfolio four times a year, the
impact cost for a fifty-thousand-dollar position works out at a
concentration of a fraction of a basis point a year. The same calculation
for a billion-dollar position in a five-hundred-million-dollar fund is
three orders of magnitude larger per unit of capital, and the square-root
law means that quadrupling the capital doubles the cost.

\begin{takeaway}[The structural advantage]
For one investor the friction is the spread; for the institution it is the
square-root law applied to a large fraction of the market's daily volume.
This is why the institutional literature's conclusions about which
strategies are implementable do not transfer, and why the single investor's
binding constraint is the quality of the forecast rather than the cost of
acting on it.
\end{takeaway}

\section{The exception: the time constraint}

There is one capacity-like constraint that a single investor has and an
institution does not, and it is worth naming because it is easy to forget.

An institution can hold a strategy for a decade while the capacity
constants change around it. A single investor who is saving for a house
in four years, or who is fifty-eight and drawing down, has a horizon that
may be shorter than the strategy's own convergence time. Chapter 8
established that the evaluation period changes the premium the investor
can harvest; this is the mirror image, and it is the reason the book's
final chapter asks the reader to state their horizon before reading any of
the allocation advice.

\begin{keytakeaways}
\item Actual institutional trading costs are less than a tenth of prior
estimates, and the value and momentum premia are the most scalable.
\item For one investor the spread is the entire friction, so strategies
that institutions cannot scale become available, and the binding
constraint becomes forecast quality rather than implementation.
\item The single investor does not escape the post-publication decay; only
the cost component of it.
\item The one capacity-like constraint the single investor has and the
institution does not is their own horizon, which may be shorter than the
strategy's convergence time.
\end{keytakeaways}
